\documentclass{ximera}

\begin{document}

    \title{Lecture 20}
    
    \begin{abstract}  
    Outline: \\
    - Properties of derivatives \\
    - Chain rule
\end{abstract}

\maketitle

\begin{definition}
    Let $g:A \to \mathbb{R}$ be a function defined on an interval $A$.  Given $c \in A$, the derivative of $g$ at $c$ is defined as
    \begin{align*}
        g'(c) := \lim\limits_{x\to c} \frac{g(x) - g(c)}{x-c},
    \end{align*}
    provided the limit exists.  If $g'$ exists for all points $c \in A$, we say that $g$ is differentiable on $A$.
\end{definition}

\begin{exercise}
    Let $f,g$ be functions defined on an interval $A$, and assume both are differentiable at some point $c \in A$.  Prove the following:
    \begin{enumerate}
        \item $(f+g)'(c) = f'(c) + g'(c)$
        \item $(fg)'(c) = f'(c)g(c) + f(c) g'(c).$
    \end{enumerate}
\end{exercise}

Now, we want to look at derivatives of quotients.  The natural way to prove this is, formally,
\begin{align*}
    &\lim\limits_{x\to c} \frac{f(x)/g(x) - f(c)/g(c)}{x-c}
    \\
    &\quad =
    \lim\limits_{x\to c} \frac{(f(x)g(c) - f(c)g(x))/(g(x)g(c))}{x-c}
    \\
    &\quad = 
    \lim\limits_{x\to c} \frac{(f(x)g(c) - f(c)g(x))}{x-c} \cdot \lim\limits_{x\to c} \frac{1}{(g(x)g(c))}
    \\
    &\quad =
    \lim\limits_{x\to c} \frac{(f(x)-f(c)g(c) - f(c)(g(x)-g(c))}{x-c} \cdot \lim\limits_{x\to c} \frac{1}{(g(x)g(c))}
    \\
    &\quad =
    \lim\limits_{x\to c} \frac{(f(x)-f(c)g(c)}{x-c} - \lim\limits_{x\to c} \frac{f(c)(g(x)-g(c))}{x-c} \cdot \lim\limits_{x\to c} \frac{1}{(g(x)g(c))}
\end{align*}

Of course, this is not correct.  Why?

\begin{example}
    Consider the function 
    \begin{align*}
        f(x) = \begin{cases}
            x^2 \sin\left(\frac1x\right) & x\neq 0 \\
            0 & x=0.
        \end{cases}       
    \end{align*}
    In this case, no matter what neighborhood I take of $x = 0$, this function evaluates to $0$ on at least one element (in fact countably many elements) in that neighborhood.  However,
    \begin{align*}
        &\lim\limits_{x\to 0} \frac{f(x) - 0}{x-0}
        \\
        &\qquad = x \sin\left(\frac1x\right) = 0,
    \end{align*}
    i.e. the derivative exists!
\end{example}
The derivative of this function is $0$ in this case and therefore the quotient rule doesn't apply, but this example is illustrative for difficulties that arise in proving chain rule, which recall from Calculus I reads as $(g\circ f)'(c) = g'(f(c))\cdot f'(c)$ (you can find similar counterexamples for the quotient rule). 
As we have now provided the formal proof for the derivative of the quotient, we will prove the chain rule (in which a similar issue arises) but now rigorously avoiding any division by $0$.

Let $g$ be an arbitrary differentiable function on $\mathbb{R}$; then we can prove chain rule in this specific case using the following naive approach (Kenton, 1999):
\begin{proof}
    Let's write out what the derivative of $g\circ f$ looks like by definition:
    \begin{align*}
        (g\circ f)'(c) &= \lim\limits_{x\to c} \frac{g(f(x)) - g(f(c))}{x-c}
        \\
        &= \lim\limits_{x\to c} \frac{g(f(x)) - g(f(c))}{f(x)-f(c)}\cdot \frac{f(x) - f(c)}{x-c}.
    \end{align*}
    Note that $f$ is continuous at $c$, so this, formally, should conclude the proof.  However, this does not hold if I have $f(x) = f(c)$ at any point in a neighborhood near $(c,f(c))$.  To overcome this problem and show the proof still holds, we instead consider the following:
    \begin{align*}
        &\sup \left( \frac{g(f(x)) - g(f(c))}{x-c} :0 < |x-c| < \delta  \right)
        \\
        &\qquad = \sup \left( \frac{g(f(x)) - g(f(c))}{x-c} :0 < |x-c| < \delta, f(x) \neq f(c) \right)
        \\
        &\qquad  = \sup \left( \frac{g(f(x)) - g(f(c))}{f(x)-f(c)}\cdot \frac{f(x) - f(c)}{x-c} :0 < |x-c| < \delta, f(x) \neq f(c) \right).
    \end{align*}
    This ensures that we are \textit{not} dividing  by $0$.  The first fraction is bounded in some neighborhood of $c$ since $g$ is assumed differentiable at $f(c)$.  The second fraction can be made arbitrarily small since $f'(c) = 0$.
\end{proof}

Now, let's prove this in generality, using a similar (but slightly more involved) trick to avoid dividing by $0$.

\begin{theorem}
    Let $f :A \to \mathbb{R}$ and $g:B\to \mathbb{R}$ satisfy $f(A) \subseteq B$ so that the composition $g\circ f$ is well-defined.  If $f$ is differentiable at $c \in A$ and if $g$ is differentiable at $f(c) \in B$, then $g \circ f$ is differentiable at $c$ with $(g\circ f)'(c) = g'(f(c))\cdot f'(c)$.
\end{theorem}
\begin{proof}
    Since $g$ is differentiable at $f(c)$, we can say
    \begin{align}
        \lim\limits_{y\to f(c)} \frac{g(y) - g(f(c))}{y-f(c)}.
    \end{align}
    This can be rewritten in the following manner: let $\tilde{d}(y)$ be the difference
    \begin{align}
        \tilde{d}(y) &= \frac{g(y) - g(f(c))}{y-f(c)} - g'(f(c)).
    \end{align}
    Clearly, $\lim\limits_{y\to f(c)} \tilde{d}(y) = 0$; consequently, define the function $d(f(c)) = 0$ and $d(y) = \tilde{d}(y)$ whenever $y \neq f(c)$.  Rearranging $d(y)$, 
    \begin{align}
        g(y) - g(f(c)) = (g'(f(c)) + d(y)) (y-f(c)).
    \end{align}
    Note this equation holds for all $y \in B$, including $y = f(c)$ (because $0=0$ is a true statement).  Thus, we are now free to substitute $y = f(t)$ for any $t \in A$.  If $t \neq c$, we can divide equation (16) by $t-c$ without penalty:
    \begin{align}
        \frac{g(f(t)) - g(f(c))}{t-c} = (g'(f(c)) + d(f(t)))\cdot \frac{f(t) - f(c)}{t-c}.
    \end{align}
    The limit can be taken now without concern of dividing by $0$.
\end{proof}


\end{document}
